\scanpage{83}
\[
\Delta\ (\subset\mathrm{Cat}_{d})\subset\mathrm{Cat}
\]
\margin{Question des générateurs des $C\in\mathrm{Ob}\,\Delta$. On ne se soucie pas que les objets généraux de $\Delta$ en aient ; mais ceux qui définissent des théories algébriques, oui\ldots}

\textcircled{1} Pour $I$ \unsure{$\mathrm{Ob}(\mathrm{Cat})$|$\in\mathrm{Ob}(\mathrm{Cat})$}\inserted{($I$ petit ?)}, le 2-foncteur
\[
\deleted{\unsure{I}}\quad\deleted{\unsure{I\in\mathrm{Cat}}}\quad C\longrightarrow\underline{\mathrm{Hom}}_{\mathrm{Cat}}(I,C)
\]
est 2-représentable (l'objet qui le 2-représente, soit $\widetilde I$, est appelé $\Delta$-enveloppe de $C$) : $\underline{\mathrm{Hom}}_{\mathrm{Cat}}(I,C)\simeq\underline{\mathrm{Hom}}_{\Delta}(\widetilde I,C)$
(i.e. l'inclusion $\Delta\longrightarrow\mathrm{Cat}$
\deleted{\textcircled{2} Pour \unsure{$C$} \unsure{$\in$}}
admet un 2-adjoint à gauche)

$I\mathbin{\unsure{-|\to|\mapsto}}\widetilde I$ est un 2 foncteur $(\mathrm{Cat})\longrightarrow\Delta$

\textcircled{2} Pour $\deleted{\unsure{B}}J\in\mathrm{Ob}\,\Delta$ et $\Sigma\subset\mathrm{Fl}(J)$\inserted{($\Sigma$ petit)}, le 2-foncteur
\[
\begin{aligned}
 C\longmapsto\underline{\mathrm{Hom}}_{\Delta}(J,C;\Sigma)
 &=\text{ss-catégorie pleine de }\underline{\mathrm{Hom}}_{\Delta}(J,C)\\
 &\quad\text{formée des }u:J\to C\text{ tels que}\\
 &\quad u(\alpha)\text{ inv. pour tt }\alpha\in\Sigma
\end{aligned}
\]
est 2-représentable (par une catégorie notée $J\Sigma_{\Delta}^{-1}$ ou $J\Sigma^{-1}$).

\uline{Conséquence}\quad Existence de 2-limites \inserted{inductives} dans $\Delta$, pour un pseudo foncteur
\[
\begin{aligned}
 I&\longrightarrow\Delta\\
 i&\longrightarrow\deleted{\unsure{J}}A_i
\end{aligned}
\]

Soit en effet $D=\varinjlim_{i,(\mathrm{Cat})}A_i$ la 2-limite inductive dans $(\mathrm{Cat})$\inserted{, avec $A_i\xrightarrow{\alpha_i}D$ dans $(\mathrm{Cat})$}, de sorte que pour tt $C\in\mathrm{Ob}(\mathrm{Cat})$ et \unsure{à|a} fortiori $C\in\mathrm{Ob}\,\Delta$
\[
\underline{\mathrm{Hom}}_{\mathrm{Cat}}(D,C)\xrightarrow{\simeq}\varprojlim\underline{\mathrm{Hom}}_{\mathrm{Cat}}(A_i,C)
\]
donc \deleted{$\underline{\mathrm{Hom}}_{\Delta}(\unsure{D},C)$}\inserted{$\underline L(C)$} ``est'' la sous-catégorie pleine de \deleted{$\underline{\mathrm{Hom}}$} $\underline{\mathrm{Hom}}_{\mathrm{Cat}}(D,C)$ \deleted{\unsure{associée au}} (décrite comme $\varprojlim\underline{\mathrm{Hom}}_{\Delta}(A_i,C)$) des $D\xrightarrow{u}C$ \inserted{dans Cat} tels que $\forall\,i$, $u\circ\alpha_i:A_i\to C$ soit dans $\underline{\mathrm{Hom}}_{\Delta}(A_i,C)$ ; \deleted{\unsure{On}} ou encore la sous-catégorie de $\underline{\mathrm{Hom}}_{\Delta}(\widetilde D,C)$ \inserted{des $\widetilde u:\widetilde D\to C$} tels que pour tt $i$, on ait $\widetilde u\circ\widetilde\alpha_i\in\mathrm{Ob}\,\underline{\mathrm{Hom}}_{\Delta}(A_i,\unsure{c|C})$\unsure{\})|)\}}, où $\widetilde\alpha_i$ est le composé $A_i\xrightarrow{\alpha_i}D\xrightarrow{\mathrm{can}}\widetilde D$. Or la condition
% archive sheet number: 82
